\chapter{Homocysteine}
\section*{What Is This Test?}
Homocysteine is a sulfur-containing amino acid formed during the metabolism of methionine, an essential amino acid. Homocysteine is not used in protein synthesis; it is remethylated to methionine (requiring folate as the methyl donor, vitamin B12 as a cofactor for methionine synthase, and vitamin B6/riboflavin) or transsulfurated to cystathionine and then cysteine (requiring vitamin B6/pyridoxal 5'-phosphate as a cofactor for cystathionine beta-synthase) \cite{selhub1999homocysteine}. Serum/plasma homocysteine is elevated (hyperhomocysteinemia) when these metabolic pathways are impaired by: (1) vitamin deficiencies (folate, vitamin B12, vitamin B6, riboflavin), (2) genetic defects in the enzymes of homocysteine metabolism (homozygous cystathionine beta-synthase [CBS] deficiency causing classic homocystinuria with severe hyperhomocysteinemia >100 mcmol/L, marfanoid habitus, ectopia lentis/downward lens dislocation, intellectual disability, and premature thromboembolism; MTHFR polymorphisms with mild hyperhomocysteinemia), (3) renal failure (impaired renal homocysteine clearance), (4) medications (methotrexate, phenytoin, nitrous oxide which inactivates vitamin B12), and (5) hypothyroidism. Homocysteine is an independent risk factor for venous thromboembolism (VTE), atherosclerotic cardiovascular disease, stroke, and peripheral arterial disease; however, randomized controlled trials have demonstrated that lowering homocysteine with folic acid, vitamin B12, and vitamin B6 supplementation does NOT reduce the risk of cardiovascular events, suggesting that hyperhomocysteinemia may be a marker rather than a causal mediator of vascular disease. Plasma homocysteine measurement is clinically indicated for: evaluation of suspected vitamin B12 or folate deficiency when clinical and hematologic findings are equivocal, diagnosis of homocystinuria in infants/children with suggestive features, and assessment of cardiovascular risk in selected patients.
\section*{Alternative Names}
Homocysteine; Plasma Homocysteine; tHcy; Total Homocysteine; Hyperhomocysteinemia Screen; Hcy
\section*{Principle}
Total homocysteine (tHcy, includes free and protein-bound homocysteine, homocystine, and mixed disulfides) is measured by immunoassay (competitive fluorescence polarization immunoassay or chemiluminescent immunoassay), HPLC with fluorescence detection, or LC-MS/MS (reference method). Specimen: fasting plasma (EDTA, lavender-top) collected on ice. Homocysteine rises 15--25\% after a methionine-rich meal (postprandial), and non-fasting samples should not be used for accurate measurement. EDTA plasma must be separated from cells within 30--60 minutes of collection because homocysteine is released from erythrocytes at room temperature, producing falsely elevated levels \cite{refsum2004facts}.
\section*{History}
Homocysteine was discovered by Vincent du Vigneaud in 1932 during studies of transmethylation, work for which he received the 1955 Nobel Prize in Chemistry. In 1962, Carson and Neill described homocystinuria in Northern Irish children with intellectual disability and ectopia lentis, and Mudd and colleagues identified cystathionine beta-synthase deficiency as the cause in 1964. In 1969, Kilmer McCully proposed the homocysteine theory of arteriosclerosis after observing premature vascular disease in a child with homocystinuria. The hypothesis was supported by the Homocysteine Studies Collaboration meta-analysis in 2002, which showed a modest, independent association between homocysteine and ischemic heart disease and stroke \cite{homocysteine2002homocysteine}. Randomized trials (HOPE-2, NORVIT, VITATOPS) subsequently showed that B-vitamin supplementation lowers homocysteine without reducing cardiovascular events, establishing that the association is not clearly causal \cite{clarke2010effects}.
\section*{Why This Test Is Ordered}
\begin{itemize}
  \item Evaluation of suspected vitamin B12 or folate deficiency when serum levels are borderline (200--300 pg/mL) and hematologic findings are equivocal
  \item Differentiation of vitamin B12 deficiency from folate deficiency when combined with serum methylmalonic acid (both are elevated in B12 deficiency; only homocysteine is elevated in isolated folate deficiency) \cite{stabler1990clinical}
  \item Diagnosis and monitoring of homocystinuria (classic CBS deficiency) in infants and children with marfanoid features, ectopia lentis, thromboembolism, or intellectual disability
  \item Investigation of unexplained venous thromboembolism or premature (age $<$50 years) arterial disease in selected patients
  \item Evaluation of patients on medications that impair homocysteine metabolism (methotrexate, phenytoin, carbamazepine, nitrous oxide)
  \item Assessment of vitamin status in chronic kidney disease, where homocysteine rises as GFR falls
\end{itemize}
\section*{Primary vs Secondary Test}
Homocysteine is a secondary, confirmatory test used to clarify equivocal vitamin B12 or folate status and to confirm suspected CBS deficiency. The primary tests in macrocytic anemia are the CBC with MCV and serum B12/folate; homocysteine and methylmalonic acid are reflex confirmatory measurements. Homocysteine is not recommended for routine cardiovascular risk screening.
\section*{How This Test Is Performed}
A fasting blood sample is collected by venipuncture into an EDTA (lavender-top) tube and placed immediately on ice. Plasma must be separated from cells within 30--60 minutes because homocysteine leaks from erythrocytes at room temperature, producing falsely elevated results. Total homocysteine (free, protein-bound, and disulfide forms) is measured by immunoassay, HPLC with fluorescence detection, or LC-MS/MS (reference method). Results are available within hours to days depending on the laboratory.
\section*{What Preparation Is Needed}
Fasting for 8--12 hours is recommended because homocysteine rises 15--25\% after a methionine-rich meal. The specimen must be chilled on ice and centrifuged promptly because homocysteine is released from red cells in whole blood at room temperature. Medications that alter homocysteine (folate, vitamin B12, methotrexate, nitrous oxide) should be documented and, when clinically feasible, held before testing.
\section*{Contraindications and Precautions}
\begin{itemize}
  \item There are no absolute contraindications. Specimen-handling precautions are critical: delays in centrifugation longer than 30--60 minutes cause spuriously elevated homocysteine. Interpretive cautions:
  \item (1) renal insufficiency elevates homocysteine independent of vitamin status;
  \item (2) MTHFR C677T homozygosity elevates homocysteine only when folate intake is low;
  \item (3) mild elevations must be interpreted with age, sex, and smoking status in mind, since homocysteine rises with age and in smokers.
\end{itemize}
\section*{Risks}
Minimal risk. Slight pain or bruising at the venipuncture site. Rarely: hematoma, infection, or vasovagal syncope.
\section*{What Happens After the Test}
Results are interpreted with serum vitamin B12, folate, methylmalonic acid, and renal function. If B12 or folate deficiency is confirmed, replacement is started and homocysteine may be rechecked after 2--4 months to document normalization. In patients with mild-to-moderate elevation due to CKD, no specific therapy is indicated, because lowering homocysteine has not been shown to reduce cardiovascular events.
\section*{Understanding Results}
\subsection*{Normal (<10 mcmol/L)}
Homocysteine metabolism intact. Adequate folate, vitamin B12, vitamin B6, and renal function. Low cardiovascular risk.
\subsection*{Mild Hyperhomocysteinemia (10--30 mcmol/L)}
Most common. Vitamin B12 or folate deficiency (the most common cause), renal insufficiency (homocysteine rises as GFR falls), MTHFR C677T homozygous (TT) genotype, hypothyroidism, medications (methotrexate, phenytoin, carbamazepine, nitrous oxide), smoking, and older age.
\subsection*{Moderate Hyperhomocysteinemia (31--100 mcmol/L)}
Severe vitamin deficiencies, moderate to severe renal failure (CKD Stage 4--5, ESRD), and heterozygosity for CBS deficiency.
\subsection*{Severe Hyperhomocysteinemia (>100 mcmol/L)}
Classic homocystinuria (homozygous CBS deficiency, the rarest cause). Presents with: marfanoid features (tall stature, long limbs, arachnodactyly), ectopia lentis (downward lens dislocation --- pathognomonic and different from Marfan syndrome's upward dislocation), intellectual disability, premature atherosclerosis and thromboembolism (arterial and venous), and psychiatric disorders. Treatment: methionine-restricted diet, betaine (remethylates homocysteine to methionine), vitamin B6, and cysteine supplementation \cite{mudd1985natural}.
\section*{Normal Results}
Fasting plasma homocysteine: 5--15 mcmol/L (adults). Optimal: <10 mcmol/L for cardiovascular risk. Moderate hyperhomocysteinemia: 15--30 mcmol/L. Intermediate: 31--100 mcmol/L. Severe: >100 mcmol/L (homocystinuria). MTHFR C677T TT genotype causes mild hyperhomocysteinemia (12--20 mcmol/L) only in the folate-deplete state; in folate-replete populations (post-fortification), homocysteine is typically normal even in TT homozygotes.
\section*{Follow-up Tests}
Serum folate (and RBC folate), serum vitamin B12, serum methylmalonic acid (elevated in B12 deficiency, distinguishing B12 from folate deficiency), serum creatinine/eGFR, TSH (hypothyroidism), plasma amino acid analysis (for homocystinuria/CBS deficiency), urinary organic acids, CBS enzyme activity, and CBS gene sequencing. MTHFR genotyping is NOT currently recommended.
\section*{References}
\bibliographystyle{unsrtnat}
\bibliography{references}
\clearpage
\section*{Regulatory Status and Authorization Date}
This chapter describes a test category, not a specific commercial product. FDA, CDSCO, EU IVDR, and MHRA approval or registration dates therefore cannot be assigned without the assay manufacturer, product name, instrument, and intended use. For laboratory-developed tests, an individual agency approval date may not exist. See the Preface for the interpretation of regulatory status.

\clearpage
